\section{Related work}
\label{sec:related}

Three bodies of work meet in this paper, and none of them measures whether verbalised
deliberation helps agents reproduce the mode shares of a real territory.

\subsection{Discrete choice and perception filters}
\label{sec:discrete-choice}

Transport research has predicted mode choice for more than fifty years with discrete choice
models, which are statistically strong and behaviourally rigid. Such a model gives each mode a
probability computed from travel time, cost and the traveller's attributes
\citep{mcfadden1974conditional,benakiva1985discrete}. Analysts fit it on travel surveys. Once
each mode has its own intercept, the fitted model reproduces the survey's mode shares on the
sample it was estimated on \citep{train2009discrete}. Its behaviour, however, is frozen at
estimation. The traveller takes the mode of highest utility, and the weights given to time and
cost, or their spread across travellers, never change afterwards.

Behavioural research documents departures that no survey records, habit among them.
Travellers with a strong mode habit acquire less information about the alternatives and use
simpler decision strategies \citep{verplanken1997habit}. Experience also filters perception.
\citet{adam2025survey} surveyed 650 respondents and found that car users rate the car as more
affordable than non-users do. The reference models we use cannot represent such
perception filters, and generative agents have been proposed to fill that gap.

\subsection{Generative agents in mobility simulation}
\label{sec:generative-mobility}

\citet{park2023generative} place a language model inside the agent, with a memory stream of
its experiences and a reflection step. \citet{liu2024toward} extend the design to travel
demand modelling, and propose a hybrid with established components as a near-term step.
CitySim \citep{bougie2025citysim} scales such agents to a city. It validates them against
time-use, travel, place-popularity and crowd-density data, but no experiment there compares a
simulated modal split with an observed one. Its successor CityReal \citep{bougie2026cityreal}
does, after tuning the textual instructions toward target population statistics while keeping
the model frozen. Its reference data are proprietary, however, and its agents are not compared
with any reference model on mode choice.

GTA (Section~\ref{sec:gap}) is the closest published system on the evaluation side, and two of
its choices limit what it can show. Each of its agents returns one mode per trip rather than a
distribution, so its scores count modes and never weigh a probability. It also simulates a
single day, and its authors name multi-day memory as future work. We add three baselines
beneath the agents and four reference models above them. Our metric scores a whole
distribution, not one mode per trip. Our agents keep events in memory across days.

\citet{alves2026evaluating} stand closest on the platform side, though their evaluation has no
human ground truth. They couple the GAMA platform to a language-model module with a persistent
memory, which decides, when a disruption occurs, whether an agent replans. They compare its
behaviour with a rule-based baseline, but never with observed behaviour.

\subsection{Distributional alignment of LLM populations}
\label{sec:alignment}

A separate literature asks whether populations of language models can stand in for human
samples. \citet{argyle2023outofone} conditioned a model on sociodemographic attributes and
recovered response patterns of the matching subpopulation. \citet{meister2024benchmarking}
found that a model verbalising a distribution aligns better than its token probabilities or
the samples it writes. Our protocol scores a verbalised distribution for that reason.

The same literature also studies these populations for what they do, not as human stand-ins.
\citet{baronchelli2025emergent} find that collective biases can emerge in decentralised model
populations even when no agent is biased individually. The SILICA instrument
\citep{bintareaf2026silica} grades claims about such populations. One of its perturbations
swaps the order in which actions are listed, and we adapt that control by randomising the
order of the options.
